\documentclass[11pt,a4paper]{article}

% ============================================================
% Relatório de CONFRONTAÇÕES TEÓRICAS — NeuraTrade
% Defesa, ponto a ponto, das decisões de modelagem do projeto, incluindo
% trade-offs e limitações assumidas. Público: banca / leitor crítico.
%
% Compilar (a partir de debate/):  tectonic main.tex
% ============================================================

\usepackage[utf8]{inputenc}
\usepackage[T1]{fontenc}
\usepackage[brazil]{babel}
\usepackage[a4paper,margin=2.5cm]{geometry}
\usepackage{amsmath,amssymb}
\usepackage{amsthm}
\usepackage{booktabs}
\usepackage{graphicx}
\usepackage{xcolor}
\usepackage{listings}
\usepackage{enumitem}
\usepackage{caption}
\usepackage{hyperref}
\usepackage{xurl}

\hypersetup{colorlinks=true,linkcolor=blue,citecolor=blue,urlcolor=blue}
\graphicspath{{../report/figures/}}

% --- Caixas de argumentação ---
\theoremstyle{definition}
\newtheorem*{confronto}{\textbf{Confrontação}}
\newtheorem*{honesto}{\textbf{Limitação assumida}}
\newtheorem*{tradeoff}{\textbf{Trade-off}}
\newtheorem*{intuicao}{Intuição}

% --- Código Python ---
\definecolor{codegray}{gray}{0.96}
\definecolor{codekw}{RGB}{0,0,170}
\definecolor{codestr}{RGB}{160,0,0}
\definecolor{codecom}{RGB}{0,120,0}
\lstdefinestyle{py}{
  language=Python, backgroundcolor=\color{codegray},
  basicstyle=\ttfamily\footnotesize, keywordstyle=\color{codekw}\bfseries,
  stringstyle=\color{codestr}, commentstyle=\color{codecom}\itshape,
  numbers=left, numberstyle=\tiny\color{gray}, numbersep=6pt,
  frame=single, framesep=4pt, rulecolor=\color{gray!40},
  breaklines=true, showstringspaces=false, tabsize=4, captionpos=b, xleftmargin=14pt,
}
\lstset{style=py}

\title{\textbf{NeuraTrade --- Confrontações Teóricas}\\[0.3em]
\large Defesa das decisões de modelagem: dados, arquitetura, validação e
estratégia de múltiplos modelos}
\author{Projeto NeuraTrade --- Redes Neurais (2026.1)}
\date{\today}

\begin{document}
\maketitle

\begin{abstract}
\noindent Este documento responde, de forma direta e crítica, às principais
objeções teóricas que o desenho do NeuraTrade pode receber. Cada seção abre com
a \textbf{Confrontação} (a pergunta difícil), apresenta a resposta fundamentada
e, quando cabe, registra explicitamente o \textbf{Trade-off} aceito e a
\textbf{Limitação assumida}. O princípio é o mesmo do projeto: preferir a
honestidade metodológica à aparência de completude. Detalhes de implementação
remetem aos \emph{Architecture Decision Records} (ADRs) em \texttt{docs/adr/};
a fundamentação didática dos conceitos está no \emph{Guia Teórico}
(\texttt{teoria/}).
\end{abstract}

\tableofcontents
\newpage

\section{Como ler este documento}
\label{sec:debate-intro}

As perguntas aqui enfrentadas são do tipo que uma banca faz para testar se o
projeto entende as \emph{consequências} de suas escolhas --- não apenas se
seguiu uma receita. Por isso cada decisão é defendida junto com o que ela
\emph{custa} e o que ela \emph{deixa de fora}.

Recapitulando o objeto em uma linha: o NeuraTrade detecta anomalias de forma
\textbf{não supervisionada} em quatro ações da B3 (PETR4, VALE3, AMER3, ITUB4),
treinando um \textbf{LSTM-Autoencoder por ativo} sobre log-retornos de
2010--2019 (``normalidade'') e sinalizando, em 2020--2024, as janelas cujo
\textbf{erro de reconstrução} ultrapassa um limiar.

\medskip
\noindent\textbf{Mapa das confrontações:}
\begin{enumerate}[noitemsep]
  \item Base de dados extremamente desbalanceada (Seção~\ref{sec:debate-desbal}).
  \item Escolha e desenho da arquitetura (Seção~\ref{sec:debate-arq}):
    por que LSTM-Autoencoder, por que os nós de entrada, profundidade e largura.
  \item Metodologia de validação (Seção~\ref{sec:debate-valid}), incluindo o
    papel da validação cruzada na seleção de modelos.
  \item Por que múltiplas redes em vez de uma rede única (Seção~\ref{sec:debate-multi}).
\end{enumerate}

\section{Base de dados extremamente desbalanceada}
\label{sec:debate-desbal}

\begin{confronto}
Anomalias são, por definição, raríssimas: talvez 1\% dos dias, ou menos. Em
classificação supervisionada, uma classe tão minoritária arruína o treino --- o
modelo aprende a dizer sempre ``normal'' e acerta 99\%. Como o projeto lida com
um desbalanceamento tão extremo?
\end{confronto}

\subsection{A resposta: o problema foi reformulado, não ``balanceado''}
A premissa da pergunta é de \emph{classificação supervisionada}. O NeuraTrade
\textbf{não} é supervisionado e \textbf{não} classifica em duas classes. A
estratégia evita o desbalanceamento na origem, reformulando a tarefa como
\textbf{aprendizado de uma classe} (\emph{one-class learning}) via reconstrução:

\begin{itemize}[noitemsep]
  \item Treina-se \textbf{somente com dados ``normais''} (2010--2019). A classe
        minoritária (anomalias) \emph{nem entra} no treino --- logo não existe
        proporção entre classes para desbalancear.
  \item O modelo aprende a \textbf{reconstruir} o normal. A anomalia é detectada
        pelo \emph{fracasso} de reconstrução (erro alto), não por um rótulo de
        classe. É a mesma lógica de detecção de fraude/falha por
        \emph{novelty detection}.
\end{itemize}

\paragraph{Intuição.} Em vez de mostrar ao modelo milhões de dias normais e um
punhado de anômalos (e implorar para ele não ignorar o punhado), mostramos
\emph{só} o normal e deixamos a anomalia ser ``tudo aquilo que ele não consegue
reproduzir''. O desequilíbrio de contagem deixa de ser uma variável do treino.

\subsection{Por que não as técnicas usuais de balanceamento}
As ferramentas clássicas contra desbalanceamento seriam contraproducentes aqui:
\begin{itemize}[noitemsep]
  \item \textbf{Oversampling / SMOTE:} sintetizaria ``anomalias'' interpolando
        exemplos --- mas não temos exemplos rotulados de anomalia, e interpolar
        em série temporal quebraria a estrutura sequencial.
  \item \textbf{Undersampling do normal:} jogaria fora justamente os dados que
        \emph{definem} a normalidade que queremos aprender.
  \item \textbf{Pesos de classe:} pressupõem duas classes no treino, que não
        existem na formulação \emph{one-class}.
\end{itemize}

\subsection{Onde o desbalanceamento reaparece --- e como o tratamos}
Ele não some de todo: ressurge na \textbf{avaliação}. Ao injetar 50 anomalias
sintéticas numa série de $\approx$1200 janelas de teste, o \emph{ground-truth} é
fortemente desbalanceado ($\approx$4\% positivos). Isso tem duas implicações que
o projeto trata explicitamente:

\begin{enumerate}[noitemsep]
  \item \textbf{Acurácia é inútil} num cenário assim (dizer ``tudo normal'' dá
        $\approx$96\% de acurácia e zero utilidade). Por isso reportamos
        \textbf{Precision, Recall e F1}, não acurácia.
  \item \textbf{Teto de precisão pela taxa-base.} O limiar p95 marca $\approx$5\%
        das janelas por construção; contra $\approx$4\% de positivos reais, há um
        piso de falsos positivos que limita a \emph{precision} máxima. Esse
        número é reportado \emph{junto} das métricas, para não inflar a leitura.
\end{enumerate}

\begin{tradeoff}
O paradigma \emph{one-class} resolve o desbalanceamento, mas paga um preço: como
nunca vemos anomalias rotuladas, não há como o modelo aprender a
\emph{discriminar} tipos de anomalia, nem otimizar diretamente uma fronteira de
decisão. Ele só sabe ``isto não parece normal''. Para um problema sem rótulos
confiáveis (o nosso caso), é a troca certa; para um cenário com rótulos
abundantes, um classificador supervisionado com ponderação de classe poderia
superar.
\end{tradeoff}

\section{Escolha e desenho da arquitetura}
\label{sec:debate-arq}

\subsection{Por que um LSTM-Autoencoder}
\label{sec:debate-arq-lstmae}

\begin{confronto}
Há dezenas de detectores de anomalia. Por que justamente um autoencoder
recorrente, e não um método estatístico clássico, um classificador, ou outra
arquitetura de rede?
\end{confronto}

A escolha resulta do cruzamento de três exigências do problema, que eliminam as
alternativas uma a uma:

\begin{enumerate}[noitemsep]
  \item \textbf{Não supervisionado} (não há rótulos). Isso descarta classificadores
        supervisionados (\emph{random forest}, redes de classificação) e métodos
        que exigem exemplos de anomalia.
  \item \textbf{Sequencial / temporal} (a ordem dos dias carrega informação; há
        dependência e regimes de volatilidade que se estendem por semanas). Isso
        enfraquece métodos que tratam cada dia como ponto independente
        (\emph{Isolation Forest}, \emph{One-Class SVM} sobre \emph{features}
        estáticas, $z$-score pontual).
  \item \textbf{Padrão complexo e não linear} de ``normalidade''. Isso desfavorece
        modelos lineares puros (ARIMA, limiar fixo sobre o retorno).
\end{enumerate}

O \textbf{autoencoder} atende (1): aprende a normalidade por reconstrução, sem
rótulos. A parte \textbf{LSTM} atende (2) e (3): processa a janela como sequência,
com memória de longo prazo e não-linearidade. A combinação --- comprimir uma
janela temporal e medir o erro ao reconstruí-la --- é o casamento natural das
três exigências.

\begin{table}[h]
\centering
\caption{Por que as alternativas foram preteridas.}
\label{tab:alternativas}
\renewcommand{\arraystretch}{1.2}
\begin{tabular}{p{4.2cm} p{9.5cm}}
\toprule
\textbf{Alternativa} & \textbf{Motivo de não adoção} \\
\midrule
$z$-score / média móvel & Linear e pontual; serve de \emph{baseline} (e foi
implementado como tal), mas não capta padrão sequencial. \\
ARIMA / GARCH & Fortes em séries financeiras, mas assumem forma paramétrica;
o objetivo do projeto é justamente uma abordagem \emph{neural} de representação aprendida. \\
Isolation Forest / OC-SVM & Não supervisionados, porém sobre vetores de
\emph{features} sem ordem temporal nativa. \\
AE com camadas densas (MLP) & Perde a estrutura sequencial: trataria os 30 passos
como 30 \emph{features} independentes. \\
AE convolucional (1D-CNN) & Capta padrão local, mas a dependência de longo alcance
é melhor servida pela memória explícita da LSTM. \\
VAE / GAN & Mais expressivos, porém mais difíceis de treinar e fora do escopo da
proposta; registrados como trabalho futuro. \\
\bottomrule
\end{tabular}
\end{table}

\begin{intuicao}
A LSTM dá ``memória'' e a estrutura de autoencoder dá ``aprender só o normal''.
Tirar qualquer uma das duas quebra um requisito: sem LSTM, ignoramos o tempo;
sem autoencoder, precisaríamos de rótulos.
\end{intuicao}

\subsection{Por que os nós de entrada (a representação de entrada)}
\label{sec:debate-arq-input}

\begin{confronto}
A entrada da rede é um tensor de forma $(30, 1)$: 30 passos no tempo, 1 variável
por passo. Por que uma única variável (e não OHLCV completo)? Por que 30 passos?
\end{confronto}

\paragraph{Uma variável: o log-retorno do fechamento.}
A entrada é o \textbf{log-retorno diário} $r_t = \ln(P_t/P_{t-1})$, e só ele.
Justificativas:
\begin{itemize}[noitemsep]
  \item \textbf{Estacionariedade.} O preço bruto tem tendência (cresce ao longo
        dos anos); o log-retorno remove essa tendência, dando à rede um alvo de
        ``normalidade'' estável no tempo. Treinar sobre preço faria o normal de
        2010 (preços baixos) parecer anômalo em 2019.
  \item \textbf{Comparabilidade entre ativos} e imunidade a \emph{splits}
        (o log transforma fatores multiplicativos em deslocamentos aditivos).
  \item \textbf{Parcimônia.} Com $\approx$2450 janelas de treino por ativo, uma
        entrada univariada mantém o número de parâmetros compatível com o volume
        de dados (Seção~\ref{sec:debate-arq-prof}), reduzindo risco de
        \emph{overfitting}.
\end{itemize}

\paragraph{Por que não OHLCV (multivariado).}
Volume e preços máximo/mínimo carregam informação real, mas: (i) multiplicariam
os parâmetros sem aumentar os dados; (ii) exigiriam normalização conjunta
cuidadosa; (iii) tornariam a detecção menos interpretável (erro misturando
escalas diferentes). A decisão é \emph{deliberadamente univariada} --- e sua
principal consequência (anomalias puramente de volume são invisíveis) é assumida
e registrada como trabalho futuro (autoencoder multivariado).

\paragraph{Por que 30 passos.}
A janela $L=30$ ($\approx$6 semanas úteis) é a granularidade da ``frase'' que a
LSTM lê. É o valor consagrado nas implementações de referência de detecção com
LSTM (Li, 2020; Valkov) e equilibra: janela curta demais não captura regime;
longa demais dilui eventos pontuais e reduz o número de janelas de treino. A
forma final $(n,\,30,\,1)$ é exatamente o que a camada LSTM espera.

\begin{tradeoff}
Univariado + janela de 30: maximiza estabilidade e parcimônia, ao custo de
cegueira a sinais fora do retorno do fechamento (volume, \emph{intraday}) e a
anomalias mais curtas que a resolução da janela.
\end{tradeoff}

\subsection{Profundidade da rede e neurônios por camada}
\label{sec:debate-arq-prof}

\begin{confronto}
Por que uma rede tão rasa? Por que 64 unidades nas LSTMs e um gargalo de 16?
Como esses números foram escolhidos --- e não chutados?
\end{confronto}

A rede tem \textbf{38.737} parâmetros, distribuídos assim:

\begin{table}[h]
\centering
\caption{Arquitetura camada a camada (entrada $(30,1)$, um modelo por ativo).}
\label{tab:camadas}
\renewcommand{\arraystretch}{1.2}
\begin{tabular}{l l r p{5.2cm}}
\toprule
\textbf{Camada} & \textbf{Saída} & \textbf{Params} & \textbf{Papel} \\
\midrule
Input (janela)            & $(30,1)$  & 0      & janela de log-retornos \\
LSTM (encoder, 64)        & $(64)$    & 16.896 & comprime a sequência em um vetor \\
Dropout (0,2)             & $(64)$    & 0      & regularização \\
Dense (gargalo, 16, ReLU) & $(16)$    & 1.040  & \emph{bottleneck}: força a compressão \\
RepeatVector(30)          & $(30,16)$ & 0      & replica o latente para o decoder \\
LSTM (decoder, 64, seq)   & $(30,64)$ & 20.736 & reconstrói passo a passo \\
Dropout (0,2)             & $(30,64)$ & 0      & regularização \\
TimeDistributed Dense(1)  & $(30,1)$  & 65     & projeta de volta ao espaço do retorno \\
\midrule
\multicolumn{2}{l}{\textbf{Total}} & \textbf{38.737} & \\
\bottomrule
\end{tabular}
\end{table}

\paragraph{Profundidade (rasa) é proporcional aos dados.}
São $\approx$2450 janelas de treino por ativo. Uma regra prática saudável é
manter os parâmetros na mesma ordem de grandeza (ou abaixo) do produto
amostras$\times$sinal, sob risco de \emph{overfitting}. Com $\approx$38,7\,mil
parâmetros para $\approx$2450 janelas de 30 passos, já estamos num regime em que
a regularização (\emph{dropout}, \emph{early stopping}) importa --- aprofundar a
rede pioraria isso sem ganho, dado que o sinal (uma série univariada) é simples.
Uma camada LSTM em cada ponta é suficiente para a expressividade necessária.

\paragraph{64 unidades: \emph{default} de literatura.}
\texttt{lstm\_units=64} e \texttt{dropout=0,2} reproduzem a implementação de
referência (Valkov), padrão consolidado em AE-LSTM para séries financeiras ---
proveniência ``\emph{default} de literatura'', não arbitrária.

\paragraph{Gargalo de 16: decisão de projeto \emph{verificada}.}
A literatura não padroniza a dimensão latente. Fixou-se 16 (compressão
$\approx 4{:}1$ sobre as 64 unidades) e, crucialmente, \textbf{mediu-se a
sensibilidade}: treinando PETR4 com $\texttt{latent\_dim}\in\{8,16,32\}$, a perda
de validação foi praticamente idêntica ($0{,}00342 / 0{,}00343 / 0{,}00342$).
Ou seja, o gargalo \emph{não é restrição ativa} nessa faixa --- até 8 bastaria,
e 16 foi mantido com folga. Isso responde diretamente à objeção do ``chute'': o
valor foi confirmado por experimento, não por fé.

\begin{honesto}
O estudo de \texttt{latent\_dim} foi conduzido em \emph{um} ativo (PETR4) e com
um único \emph{seed}. A planura observada é forte evidência, mas uma varredura
nos quatro ativos e com múltiplos \emph{seeds} seria mais robusta. Igualmente,
\texttt{lstm\_units} e \texttt{dropout} foram herdados da literatura, sem
varredura própria --- assumido como escopo.
\end{honesto}

\section{Metodologia de validação}
\label{sec:debate-valid}

\begin{confronto}
Como o projeto valida os modelos? Há um conjunto de teste honesto? A separação
respeita a natureza temporal dos dados, ou há vazamento?
\end{confronto}

\subsection{Três níveis de separação}
A validação opera em três níveis, deliberadamente distintos:

\begin{enumerate}[noitemsep]
  \item \textbf{Treino} (2010--2019): aprende a normalidade.
  \item \textbf{Validação} (bloco final do treino, $\approx$2018--2019, via
        \texttt{validation\_split=0,1} com \texttt{shuffle=False}): governa o
        \emph{early stopping} e a calibração de hiperparâmetros. \textbf{Não} é o
        teste.
  \item \textbf{Teste} (2020--2024): nunca visto no treino; é onde as anomalias
        são detectadas e as métricas reportadas.
\end{enumerate}

\paragraph{O detalhe crítico: \texttt{shuffle=False}.}
O \texttt{validation\_split} do Keras separa a \emph{última fração} do array
\emph{antes} de embaralhar. Com \texttt{shuffle=True}, cada época re-embaralharia
o treino e --- combinada com janelas sobrepostas --- misturaria amostras
temporalmente adjacentes entre treino e validação, vazando padrão futuro. Com
\texttt{shuffle=False}, treino e validação ficam em blocos cronológicos
contíguos. Esta é a defesa central contra vazamento, e está registrada como
decisão crítica (ADR-0004).

\begin{honesto}
Há um vazamento \emph{residual} conhecido: janelas na fronteira treino/validação
compartilham alguns dias, por causa da sobreposição de janelas. O efeito é
pequeno e foi assumido e documentado, não ignorado.
\end{honesto}

\subsection{Validação cruzada para seleção de modelos}
\label{sec:debate-valid-cv}

\begin{confronto}
Para escolher o modelo ótimo entre opções pré-determinadas, usou-se
\emph{validação cruzada}? Onde está o $k$-fold?
\end{confronto}

Aqui é preciso ser direto: \textbf{não usamos $k$-fold clássico --- e isso é
intencional, não um esquecimento.}

\paragraph{Por que $k$-fold padrão é inválido aqui.}
A validação cruzada $k$-fold embaralha e particiona os dados assumindo que as
amostras são \emph{independentes e identicamente distribuídas} (i.i.d.). Séries
temporais violam isso: as observações são dependentes e ordenadas. Um \emph{fold}
de treino contendo dias \emph{posteriores} aos do \emph{fold} de teste constitui
vazamento de futuro --- exatamente o pecado da Seção~\ref{sec:debate-valid}. Em
detecção de anomalias isso é ainda pior, pois infla artificialmente a aparência
de desempenho.

\paragraph{O que de fato fizemos: seleção por \emph{grid} sobre validação temporal.}
A ``comparação de múltiplos modelos para escolher o ótimo'' aconteceu, mas pela
via apropriada a séries temporais --- avaliar candidatos numa validação que
respeita o tempo:
\begin{itemize}[noitemsep]
  \item \textbf{Dimensão latente} $\in\{8,16,32\}$: comparadas pela
        \texttt{val\_loss} no bloco de validação cronológico; resultado plano,
        16 confirmado (Seção~\ref{sec:debate-arq-prof}).
  \item \textbf{Esquema de limiar} (estático p95 vs.\ dinâmico causal) e
        \textbf{tamanho da janela dinâmica} $\in\{126,252,504\}$: comparados pela
        taxa de marcação e pelo comportamento sob mudança de regime no teste; 252
        confirmado.
  \item \textbf{Modelo profundo vs.\ \emph{baseline}}: o LSTM-AE é confrontado
        com um \emph{baseline} de $z$-score em janela móvel, para evidenciar
        ganho sobre a heurística trivial.
\end{itemize}

\paragraph{A forma correta de CV temporal (e nosso caminho de rigor).}
A validação cruzada que \emph{seria} válida aqui é a de série temporal ---
\emph{walk-forward} / janela expansível, em que cada \emph{fold} treina no
passado e valida no futuro imediato, nunca o contrário:

\begin{lstlisting}[caption={Validação cruzada temporal (walk-forward) com scikit-learn.}]
from sklearn.model_selection import TimeSeriesSplit
tscv = TimeSeriesSplit(n_splits=5)
for treino_idx, val_idx in tscv.split(X):
    # treino_idx sempre ANTES de val_idx -> sem vazamento de futuro
    modelo = build_lstm_autoencoder()
    modelo.fit(X[treino_idx], X[treino_idx], shuffle=False, ...)
    avalia(modelo, X[val_idx])
\end{lstlisting}

\begin{honesto}
O projeto usou um \textbf{único} \emph{holdout} temporal de validação (não uma
$k$-fold temporal com vários cortes). Para a seleção de hiperparâmetros conduzida
--- de baixa dimensionalidade e com resultados planos/robustos --- isso é
suficiente e defensável. Porém, um \texttt{TimeSeriesSplit} (\emph{walk-forward})
daria estimativas com \emph{variância} (média $\pm$ desvio entre cortes) e uma
seleção de modelo estatisticamente mais sólida. Está registrado como o próximo
passo de rigor --- e o ADR-0004 já o aponta como alternativa considerada.
\end{honesto}

\subsection{Validação substantiva: além da perda}
Vale frisar que a \texttt{val\_loss} não é o único juiz. Como não há rótulos de
anomalia reais, a validação do \emph{detector} (não só do reconstrutor) se dá por
duas vias complementares (detalhadas no relatório principal):
\begin{itemize}[noitemsep]
  \item \textbf{Quantitativa:} injeção sintética de anomalias em posições
        conhecidas $\to$ Precision/Recall/F1.
  \item \textbf{Narrativa:} as anomalias detectadas coincidem com eventos reais
        (COVID, fraude da Americanas, Brumadinho)?
\end{itemize}
Essa dupla checagem é a resposta honesta para ``como validar um modelo sem
gabarito''.

\section{Múltiplas redes vs.\ uma rede única}
\label{sec:debate-multi}

\begin{confronto}
Treinar quatro modelos separados (um por ativo) desperdiça dados e esforço. Por
que não uma única rede, treinada com os quatro ativos juntos, que generalizaria
melhor?
\end{confronto}

\subsection{A razão de fundo: ``normalidade'' é específica do ativo}
O detector funciona aprendendo \emph{o que é normal}. Mas a normalidade de cada
ativo é distinta: a volatilidade típica, a autocorrelação e o regime de risco de
PETR4 (energia, sensível a petróleo e política) não são os de ITUB4 (banco,
mais estável) nem os de AMER3 (varejo, que entrou em colapso). Uma rede única
aprenderia uma normalidade \emph{média} --- e a média de comportamentos
heterogêneos não descreve bem nenhum deles.

\begin{intuicao}
Um médico que atende só cardiologistas aprende com precisão o ``normal'' daquele
grupo. Um clínico que mistura recém-nascidos, atletas e idosos num só padrão de
referência classificaria mal todos: o batimento normal de um bebê alarmaria como
anômalo num idoso. Cada ativo é uma ``população'' com sua fisiologia.
\end{intuicao}

\subsection{O objetivo do projeto exige a separação}
A proposta do NeuraTrade é \textbf{comparar a generalização entre setores}
(energia, mineração, varejo, financeiro). Essa comparação \emph{só faz sentido}
com um modelo por ativo: é o que permite afirmar ``o detector de AMER3 cobre 9
de 10 dos seus eventos'' ou ``um choque sistêmico (COVID) acende em todos os
quatro''. Uma rede única dissolveria a unidade de análise (o ativo/setor) que é
o próprio objeto de estudo.

\subsection{Os contra-argumentos, honestamente}
A rede única tem vantagens reais, que reconhecemos:

\begin{table}[h]
\centering
\caption{Modelo por ativo (adotado) vs.\ rede única (alternativa).}
\label{tab:multi}
\renewcommand{\arraystretch}{1.2}
\begin{tabular}{p{4.0cm} p{4.6cm} p{4.9cm}}
\toprule
 & \textbf{Um modelo por ativo (adotado)} & \textbf{Rede única (todos os ativos)} \\
\midrule
Especificidade & Alta: capta o regime de cada ativo & Baixa: aprende uma média \\
Dados por modelo & $\approx$2450 janelas & $\approx$9800 janelas (4$\times$) \\
Comparação setorial & Direta e interpretável & Indireta / perdida \\
Risco de \emph{overfitting} & Maior (menos dados) & Menor (mais dados) \\
Custo / manutenção & 4 modelos para treinar e versionar & 1 modelo \\
Transferência entre ativos & Nenhuma & Implícita (pode ajudar ativos com pouco histórico) \\
\bottomrule
\end{tabular}
\end{table}

\paragraph{Por que, no saldo, a separação vence aqui.} O volume de dados por
ativo ($\approx$2450 janelas) é \emph{suficiente} para a rede pequena adotada
(Seção~\ref{sec:debate-arq-prof}), então o principal trunfo da rede única (mais
dados) não é decisivo. Já a especificidade da normalidade e a exigência de
comparação setorial são centrais ao problema. O custo de manter quatro modelos
pequenos é baixo e a configuração é centralizada (um \texttt{config.yaml}, o
mesmo código, só o ativo muda).

\begin{tradeoff}
Modelo por ativo troca ``mais dados e um artefato único'' por ``especificidade e
comparabilidade setorial''. Para um problema cujo \emph{objetivo declarado} é
comparar setores, a troca é favorável. Num cenário com centenas de ativos, ou
com ativos de histórico curto, uma \textbf{arquitetura intermediária} ---
um modelo global com \emph{embedding} de ativo, ou \emph{fine-tuning} por ativo
sobre uma base compartilhada --- seria o próximo passo natural.
\end{tradeoff}

\section{Síntese e limitações assumidas}
\label{sec:debate-sintese}

As confrontações compartilham um fio condutor: cada decisão do NeuraTrade
otimiza para o \emph{problema real} --- detecção não supervisionada, em séries
temporais, com normalidade específica por ativo --- e não para um \emph{checklist}
genérico de boas práticas. Resumindo as defesas:

\begin{itemize}[noitemsep]
  \item \textbf{Desbalanceamento:} dissolvido pela reformulação \emph{one-class}
        (treina só no normal); reaparece na avaliação e é tratado com P/R/F1 e
        reporte da taxa-base, nunca acurácia.
  \item \textbf{Arquitetura (LSTM-AE):} única combinação que satisfaz
        simultaneamente ``sem rótulos'', ``sequencial'' e ``não linear''. Entrada
        univariada e janela de 30 por estacionariedade e parcimônia. Rede rasa e
        gargalo 16 dimensionados ao volume de dados, com a dimensão latente
        \emph{verificada} por sensibilidade.
  \item \textbf{Validação:} três níveis com separação temporal estrita
        (\texttt{shuffle=False}); seleção de modelos por \emph{grid} sobre
        validação temporal; $k$-fold clássico evitado por ser inválido em séries
        temporais.
  \item \textbf{Múltiplas redes:} a normalidade é específica do ativo e a
        comparação setorial é o objetivo --- ambos exigem um modelo por ativo.
\end{itemize}

\subsection{O que ficou em aberto (e por quê)}
Listamos junto, para não dispersar, as limitações assumidas ao longo do texto:

\begin{enumerate}[noitemsep]
  \item \textbf{Validação cruzada temporal completa} (\texttt{TimeSeriesSplit} /
        \emph{walk-forward}) em vez do \emph{holdout} único, para estimativas com
        variância (Seção~\ref{sec:debate-valid-cv}).
  \item \textbf{Sensibilidade de hiperparâmetros mais ampla:} \texttt{latent\_dim}
        nos quatro ativos e múltiplos \emph{seeds}; varredura de
        \texttt{lstm\_units} e \texttt{dropout} (Seção~\ref{sec:debate-arq-prof}).
  \item \textbf{Entrada multivariada (OHLCV):} habilitaria detectar anomalias de
        volume, hoje invisíveis (Seção~\ref{sec:debate-arq-input}).
  \item \textbf{Erro por máximo na janela} (em vez da média), para combater a
        diluição de choques pontuais e elevar o \emph{recall}.
  \item \textbf{Arquitetura intermediária} entre ``um por ativo'' e ``rede única''
        (modelo global com \emph{embedding} de ativo), relevante ao escalar para
        muitos ativos (Seção~\ref{sec:debate-multi}).
\end{enumerate}

Nenhum desses itens invalida os resultados atuais; todos são caminhos de
fortalecimento. O compromisso do projeto é declarar essas fronteiras
explicitamente --- a alternativa (silenciá-las) seria a única falha
metodológica de fato grave.

\section{Da limitação à agenda: o que a Fase 2 corrige}
\label{sec:debate-fase2}

\begin{confronto}
``Vocês listaram com honestidade as fraquezas (Seção~\ref{sec:debate-sintese}), mas listar não é
corrigir. Sem um plano concreto, a transparência vira álibi para a inação.''
\end{confronto}

A crítica é justa e foi antecipada. A segunda fase do projeto (milestone~M8) converte três das
fronteiras assumidas em tarefas com decisão registrada (ADR-0009, 0010 e 0011, em status
\emph{proposto}) e ordem de prioridade explícita. O ponto desta seção não é apresentar
resultados --- não há ainda ---, mas mostrar que cada limitação tem um caminho de correção
\emph{já desenhado}, com os riscos mapeados.

\subsection{Diluição do choque $\to$ agregação por máximo/percentil}
\label{sec:debate-fase2-pooling}
A média sobre 30 passos atenua um choque de um dia (Seção~\ref{sec:debate-sintese}, item~4). A
correção troca a agregação por \emph{max} ou percentil alto, recuperando o pico.
\begin{tradeoff}
Maior \emph{recall} de choques pontuais ao custo de potencial perda de \emph{precision}: o máximo
é sensível também ao ruído de um único dia. Mitiga-se com o percentil (p90) em vez do máximo puro,
e mede-se o efeito real com P/R/F1 nos quatro ativos.
\end{tradeoff}
\begin{honesto}
A mudança invalida o limiar p95 atual (calibrado sobre a média): ele \textbf{tem} de ser
recalibrado sobre o novo escore. Reportar Recall de \emph{max} com limiar de \emph{mean} seria um
erro de comparação --- e está explicitamente vedado no plano (ADR-0009).
\end{honesto}

\subsection{\emph{Holdout} único $\to$ validação \emph{walk-forward}}
\label{sec:debate-fase2-cv}
O item~1 da síntese --- validação cruzada temporal --- materializa-se com \texttt{TimeSeriesSplit}
(ADR-0010), respondendo à objeção de variância levantada na Seção~\ref{sec:debate-valid-cv}.
\begin{honesto}
A validação cruzada não melhora o detector; melhora a \emph{confiança} na escolha de
hiperparâmetros. É instrumento de credibilidade, não de desempenho --- e é assim que será
reportada, sem inflar sua contribuição.
\end{honesto}
A regra anti-vazamento (ADR-0001) é preservada \textbf{por fold}: scaler reajustado só no treino
de cada corte, janelas geradas dentro da partição.

\subsection{Univariado $\to$ entrada multivariada (OHLCV)}
\label{sec:debate-fase2-multi}
Os itens~2 e~3 da síntese (anomalias de volume, hoje invisíveis) tornam-se a transição para um
tensor multivariado (ADR-0011), começando por Close+Volume.
\begin{tradeoff}
Ganha-se um canal precursor (volume) ao custo de $5\times$ a entrada, mais capacidade e mais
risco de \emph{overfitting} --- razão pela qual depende da validação \emph{walk-forward} já
disponível e começa pelo mínimo informativo, não pelo OHLCV completo.
\end{tradeoff}

\subsection{Por que nesta ordem}
\label{sec:debate-fase2-ordem}
A prioridade não é arbitrária: \textbf{(1)} a agregação é o ganho rápido de \emph{recall} que
custa poucas linhas; \textbf{(2)} a validação \emph{walk-forward} é a régua honesta para
comparar (1) contra a \emph{baseline} e calibrar (3); \textbf{(3)} o multivariado é a aposta de
maior alcance, viável apenas depois que o sinal deixa de ser diluído e os números passam a ser
confiáveis. Corrigir na ordem inversa seria medir uma aposta cara com um instrumento que ainda
esconde o que se quer medir.

\section{Contexto macro: por que condicionar, e não empilhar}
\label{sec:debate-macro}

\begin{confronto}
``Adicionar Selic, IPCA, dólar e VIX ao tensor de entrada é o caminho óbvio para dar contexto
sistêmico. Por que inventar um \emph{Conditional} Autoencoder em vez de simplesmente passar de
$(30,2)$ para $(30,2{+}N)$ e reconstruir tudo?''
\end{confronto}

Porque ``empilhar e reconstruir tudo'' \textbf{não atinge o objetivo} --- e isso é demonstrável,
não opinião. Três falhas:

\begin{enumerate}[noitemsep]
  \item \textbf{O MSE ignora a macro.} Indicadores de baixa frequência (IPCA mensal, Selic em
        escada), após o preenchimento causal, são quase constantes na janela. Sua variância escalada
        é ínfima; reconstruí-los é trivial; o gradiente vem só de preço/volume. A macro entraria como
        \emph{decoração inerte}.
  \item \textbf{Eventos agendados viram falsos positivos.} Uma alta de Selic ou um \emph{print} de
        IPCA são \emph{programados}. Se a macro está no alvo e o erro é agregado por \emph{max}
        (Seção~\ref{sec:debate-fase2-pooling}), um evento de calendário dispararia uma ``anomalia''.
  \item \textbf{Um escalar único apaga a distinção.} Colapsar preço/volume e macro num só erro
        impede dizer \emph{qual} bloco se mexeu --- justamente a informação que separa idiossincrático
        de sistêmico.
\end{enumerate}

A solução é o \textbf{Conditional Autoencoder}: o encoder \emph{vê} a macro (ela condiciona a
codificação), mas o decoder reconstrói \textbf{só} preço/volume e a perda recai só sobre esse bloco.
O estresse macro é medido \emph{direto da entrada} (amplitude na janela), em um eixo separado do
erro do ativo.

\begin{tradeoff}
Ganha-se a distinção idiossincrático \emph{vs.} sistêmico ao custo de uma arquitetura assimétrica
(entrada $\ne$ alvo) e de um conector de dados macro com tratamento de frequências mistas. A prova
de conceito (notebook \texttt{12}) pagou: COVID/2020 saiu \textbf{sistêmico}; Americanas/2023,
\textbf{idiossincrático}.
\end{tradeoff}

\begin{honesto}
A defesa anti-vazamento depende de um \textbf{contrato}: a macro é indexada pela \emph{data de
publicação}, não de referência (o IPCA de março só existe para a rede em meados de abril), com
\emph{ffill} causal e sem \emph{backfill}. Assumimos um \emph{lag} de publicação \textbf{aproximado}
para o IPCA --- não a data exata de divulgação ---, o que é uma simplificação a refinar. E as séries
mensais (Selic, IPCA) acabam quase inertes na janela: o contexto sistêmico vem, na prática, dos
indicadores diários (USD/BRL, VIX). O método funciona, mas o crédito honesto vai para os dois canais
diários, não para os quatro.
\end{honesto}

\begin{honesto}
Esta extensão (ADR-0012) é uma \textbf{prova de conceito} --- ainda não incorporada como entrega
oficial. O \texttt{latent\_dim} do espaço ampliado carece de revalidação por \emph{walk-forward}
antes de qualquer número final.
\end{honesto}


\end{document}
